\problemname{Från gem till hus}

Vi har alla hört den fascinerande berättelsen om hur någon lyckades byta sig från ett enda gem 
hela vägen till ett hus genom en serie smarta byten. 
I dagens bostadsmarknad låter detta som en lysande idé – men nu visar det sig att 
ett enda gem inte räcker. Din uppgift är att ta reda på det \emph{minsta antal gem} som krävs 
för att i slutändan kunna få ett hus.

Du får en uppsättning möjliga byten, där varje byte har formen
\[
3 \cdot \text{item}_1 + 2 \cdot \text{item}_2 + \dots = 1 \cdot \text{item}_k.
\]
Detta innebär att du betalar en linjärkombination av olika föremål (items), 
och i utbyte får du exakt \emph{ett} styck av \(\text{item}_k\). 
Varje byte kan utföras valfritt antal gånger.

\section*{Indata}
\begin{itemize}
    \item Låt \(N\) vara det totala antalet termer i alla byten (summan av alla ingående 
          items i vänsterledet över samtliga byten).
    \item Det gäller att \(2 \le N \le 10^5\).
    \item Koefficienterna (talen framför respektive item i bytena) är mindre än eller 
          lika med 100.
    \item Du kan använda varje byte hur många gånger du vill.
\end{itemize}

\section*{Utdata}
Skriv ut ett heltal: det minsta antalet \emph{gem} som krävs för att du ska kunna 
byta dig till ett \emph{hus}. 
Det är garanterat att svaret inte överstiger \(10^9\).

\section*{Delpoäng (Subtasks)}
\noindent
\begin{tabular}{| l | l | p{10cm} |}
  \hline
  \textbf{Grupp} & \textbf{Poäng} & \textbf{Gränser} \\ \hline
  1 & ? & Endast tre saker: \emph{gem}, \emph{hus} och \emph{ost}. \\ \hline
  2 & ? & Varje linjärkombination har exakt en term. \\ \hline
  3 & ? & \(N \leq 1000\). \\ \hline
  4 & ? & Alla koefficienter i linjärkombinationerna är 1 eller 0. \\ \hline
  5 & ? & Inga ytterligare begränsningar. \\ \hline
\end{tabular}
